% Provenance (never printed): see sources/outline_sources.md (Year 2: eccentricity, obliquity and
% precession; insolation computed from spherical astronomy; the Quaternary cycles; carbon dioxide
% and feedbacks). Orbital elements and tabulated insolation of Berger and Loutre (1991) via NOAA
% NCEI (data/derived/bachelor-2-orbit91, fetch_data Book 4 block); periods after Laskar et al.
% (2004) and the NASA Milankovitch page; history after the NASA Earth Observatory feature (2000);
% terminations after Barker and Knorr (2021); the mid-Pleistocene transition after Clark et al.
% (2006) and Koehler and van de Wal (2020). Figures from figdata/bachelor-2/orbital.py; values of
% the exercises and problem computed in the Book 4 scratch calculator calc/ch25.py.
\chapter{O forçamento orbital e as eras glaciais}\label{ch:b2:orbital-forcing}

No último milhão de anos a forma da órbita terrestre alterou a luz solar recebida pelo
planeta inteiro num ano em menos de duas partes por mil: entre \num{340.00} e
\qty{340.54}{W.m^{-2}} no topo da atmosfera. No mesmo milhão de anos a insolação média diária
de meados de julho a \ang{65} N oscilou entre cerca de 388 e \qty{500}{W.m^{-2}}, mais de
um quarto. As eras glaciais seguiram o segundo número, não o primeiro. Este capítulo calcula ambos
a partir dos elementos da órbita terrestre, e pergunta como a luz solar numa latitude, numa estação, pode
ritmar o crescimento e o recuo de mantos de gelo que cobrem continentes.
% ledger: data:b4-orbital, berger91:readme

\begin{recall}
\textbf{Dos livros de física}: a elipse, seu semieixo maior $a$ e sua excentricidade $e$, as leis de
Kepler. \textbf{Do volume do primeiro ano}: a eclíptica e a obliquidade $\varepsilon$, a declinação
do Sol $\delta$ com $\sin\delta = \sin\varepsilon\sin\lambda$ ($\lambda$ a longitude verdadeira do Sol,
contada ao longo da eclíptica a partir do equinócio vernal), o ano trópico, a insolação média
diária
\[
  \bar Q = \frac{S_0}{\pi}\left(\frac{a}{r}\right)^{2}\bigl(h_0\sin\varphi\sin\delta +
  \cos\varphi\cos\delta\sin h_0\bigr),\qquad \cos h_0 = -\tan\varphi\tan\delta ,
\]
e sua média global e anual, $S_0/(4\sqrt{1 - e^2})$. \textbf{Dos volumes escolares}: o
ciclo glacial e o Último Máximo Glacial; o forçamento radiativo. \textbf{De capítulos anteriores}: a
\omterm{def:b2:precambrian-earth:ice-albedo}{retroalimentação gelo--albedo} (\cref{def:b2:precambrian-earth:ice-albedo}); a compilação bentônica, o registro de EPICA
Dome C e os \omterm{def:b2:climate-proxies:age-model}{modelos de idade} (\cref{ch:b2:climate-proxies}).
\end{recall}
% ledger: earth:TSI, earth:obliquity, earth:ecc

\section{A órbita e seus três elementos}\label{sec:b2:orbital-forcing:elements}

A órbita terrestre é uma elipse com o Sol num dos focos. Os planetas, Júpiter e Saturno acima de
tudo, deformam-na lentamente e fazem-na girar em seu plano; a atração do Sol e da Lua sobre o bojo
equatorial faz o eixo de rotação oscilar como um pião. Três grandezas descrevem o que importa para a
luz solar: quão alongada é a elipse, quão inclinado é o eixo, e em que ponto da órbita caem as
estações.
% ledger: nasa:milank-causes

\begin{definition}[Longitude do periélio; precessão climática]\label{def:b2:orbital-forcing:elements}
A \emph{longitude do periélio}\index{longitude do periélio} $\varpi$ é o ângulo, medido no
plano da órbita a partir da direção do equinócio vernal, entre essa direção e a
direção do periélio vista do Sol; ela é contada a partir do equinócio móvel, de modo que
varia tanto com a oscilação do eixo quanto com a rotação da órbita. A \emph{precessão
climática}\index{precessão climática} é a grandeza $e\sin\varpi$, também chamada índice de
precessão.
\end{definition}
% ledger: berger91:readme

\begin{omfigure}
\centering
\begin{tikzpicture}[scale=0.82, font=\scriptsize, >=Stealth]
\draw[thick, rotate=101.37] (-0.72,0) ellipse (2.4 and 2.2895);
\draw[densely dashed, black!60] (-0.331,1.647) -- (0.615,-3.059);
\draw[->, black!70] (0,0) -- (3.9,0);
\node[anchor=south, align=center] at (3.75,0.06) {$\gamma$, em direção ao\\equinócio vernal};
\draw[->, omDef] (0.75,0) arc[start angle=0, end angle=101.37, radius=0.75];
\node[omDef] at (0.78,0.62) {$\varpi$};
\fill[ocRed] (0,0) circle (0.13);
\node[anchor=north west] at (0.05,-0.08) {Sol};
\foreach \x/\y in {-2.062/0, 0/-3.094, 2.321/0, 0/1.688} {\fill[ocBlue] (\x,\y) circle (0.08);}
\node[anchor=east, align=right] at (-2.15,0) {equinócio\\de março};
\node[anchor=north west, align=left] at (2.38,-0.08) {equinócio\\de setembro};
\node[anchor=west] at (0.1,1.92) {solstício de dezembro};
\node[anchor=east] at (-0.12,-3.25) {solstício de junho};
\fill[black] (-0.331,1.647) circle (0.06);
\node[anchor=east] at (-0.42,1.72) {periélio};
\fill[black] (0.615,-3.059) circle (0.06);
\node[anchor=west] at (0.7,-3.15) {afélio};
\draw[->, thick, black!60] (-1.9,-1.75) arc[start angle=220, end angle=245, radius=2.6];
\end{tikzpicture}
\par\smallskip{\small A órbita terrestre vista do lado norte, com a excentricidade exagerada
para 0.3 e a \omterm{def:b2:orbital-forcing:elements}{longitude do periélio} atual, $\varpi = \ang{101.4}$. A Terra move-se no sentido anti-horário
(seta); os pontos marcam sua posição nos equinócios e solstícios. O periélio cai alguns
dias depois do solstício de dezembro, o afélio alguns dias depois do solstício de junho. Esquemático.}
\end{omfigure}
% ledger: berger91:rows, nasa:milank-detail

Vista da Terra, o Sol percorre a eclíptica, e sua longitude verdadeira $\lambda$ é o
ângulo a partir do equinócio vernal: $\lambda = \ang{90}$ no solstício de junho, \ang{270} no solstício de
dezembro. A Terra está no periélio quando o Sol está na direção oposta à do periélio,
em $\lambda = \varpi + \ang{180}$; hoje $\varpi = \ang{101.37}$ e o periélio cai em
$\lambda = \ang{281.37}$, cerca de onze dias depois do solstício de dezembro, por volta de 3 de janeiro.
% ledger: berger91:rows, nasa:milank-detail

\begin{definition}[Parâmetros de Milankovitch]\label{def:b2:orbital-forcing:milankovitch}
Os \emph{parâmetros de Milankovitch}\index{parâmetros de Milankovitch} são as três grandezas
orbitais que fixam a distribuição da luz solar sobre a Terra ao longo das estações e das
latitudes: a excentricidade $e$, a obliquidade $\varepsilon$ e a \omterm{def:b2:orbital-forcing:elements}{precessão climática}
$e\sin\varpi$. A \emph{teoria de Milankovitch}\index{teoria de Milankovitch} afirma que suas lentas
variações ritmam o crescimento e o recuo dos mantos de gelo quaternários, através da insolação de verão
nas altas latitudes do norte.
\end{definition}
% ledger: nasa:milank-detail, eo00:milankovitch

Cada parâmetro varia com seus próprios ritmos. A excentricidade tem períodos principais perto de \num{95000},
\num{99000}, \num{124000} e \num{131000} anos, que se combinam num ciclo de cerca de
\num{100000} anos, e um período de \num{405000} anos, o mais estável de todos; no último
milhão de anos ela variou de \num{0.0046} a \num{0.0565}, e hoje, em \num{0.0172}, diminui
lentamente. A obliquidade varia com um período principal de \num{40978} anos, entre \ang{22.1} e
\ang{24.5} no último milhão de anos; em \ang{23.44}, ela diminui, passado seu último máximo
há cerca de \num{10000} anos. O eixo dá uma volta completa em cerca de \num{25770} anos, enquanto a
órbita gira em cerca de \num{112000}; juntos, deslocam o periélio pelas estações com
períodos principais perto de \num{19000}, \num{22400} e \num{23700} anos.
% ledger: la04:eccentricity, la04:obliquity, la04:precession, nasa:milank-detail, data:b4-orbital

\begin{omfigure}
\centering
% figdata/bachelor-2/orbital.py part elements
\begin{tikzpicture}
\begin{groupplot}[group style={group size=1 by 3, vertical sep=0.3cm, xlabels at=edge bottom,
  xticklabels at=edge bottom}, width=0.96\linewidth, height=3.1cm, xmin=0, xmax=1000,
  xtick={0,100,...,1000}, tick label style={font=\scriptsize}, label style={font=\small},
  scaled y ticks=false, yticklabel style={/pgf/number format/fixed,
  /pgf/number format/precision=2}]
\nextgroupplot[ylabel={$e$}, ymin=0, ymax=0.085, ytick={0,0.02,0.04,0.06}]
\addplot[omDef, thin] table[x=age_ka, y=ecc] {figdata/out/bachelor-2-orbital-elements.dat};
\node[font=\scriptsize, anchor=north west] at (rel axis cs:0.01,0.97) {a\enspace excentricidade};
\nextgroupplot[ylabel={$\varepsilon$ (\unit{\degree})}, ymin=21.9, ymax=25.4, ytick={22,23,24}]
\addplot[ocBlue, thin] table[x=age_ka, y=obl] {figdata/out/bachelor-2-orbital-elements.dat};
\node[font=\scriptsize, anchor=north west] at (rel axis cs:0.01,0.97) {b\enspace obliquidade};
\nextgroupplot[ylabel={$e\sin\varpi$}, ymin=-0.07, ymax=0.12, ytick={-0.05,0,0.05},
  xlabel={idade (milhares de anos antes de 1950)}]
\addplot[ocRed, thin] table[x=age_ka, y=prec] {figdata/out/bachelor-2-orbital-elements.dat};
\node[font=\scriptsize, anchor=north west] at (rel axis cs:0.01,0.97) {c\enspace precessão
climática};
\end{groupplot}
\end{tikzpicture}
\par\smallskip{\small Os três \omterm{def:b2:orbital-forcing:milankovitch}{parâmetros de Milankovitch} no último milhão de anos: (a) a
excentricidade; (b) a obliquidade; (c) a \omterm{def:b2:orbital-forcing:elements}{precessão climática}, cuja envoltória é a excentricidade.
Dados: Berger e Loutre (1991), NOAA NCEI Paleoclimatology, obtidos em 23 de outubro de 2026.}
\end{omfigure}
% ledger: data:b4-orbit91

A \omterm{def:b2:orbital-forcing:elements}{precessão climática} oscila depressa, mas só até onde a excentricidade permite: quando a órbita é
quase circular, como há \num{400000} anos, pouco importa onde cai o periélio. A
figura mostra-o claramente: a envoltória do painel (c) é o painel (a).
% ledger: data:b4-orbit91

\section{A luz solar a partir da órbita}\label{sec:b2:orbital-forcing:insolation}

A insolação média diária lembrada acima precisa de mais um ingrediente: a distância do Sol num
dado dia.

\begin{proposition}[A distância e a longitude do Sol]\label{prop:b2:orbital-forcing:distance}
Quando a longitude verdadeira do Sol é $\lambda$, a distância Terra--Sol é
\[
  \frac{r}{a} = \frac{1 - e^2}{1 - e\cos(\lambda - \varpi)} .
\]
No solstício de junho ($\lambda = \ang{90}$), em primeira ordem em $e$,
$(a/r)^2 \approx 1 - 2e\sin\varpi$; no solstício de dezembro, $(a/r)^2 \approx 1 + 2e\sin\varpi$.
\end{proposition}

\begin{proof}
Numa elipse com o Sol num foco, $r/a = (1 - e^2)/(1 + e\cos v)$, onde a anomalia verdadeira $v$
é o ângulo no Sol do periélio à Terra. A Terra é vista do Sol na
direção oposta à do Sol visto da Terra, de modo que sua longitude vista do Sol é
$\lambda + \ang{180}$, e $v = \lambda + \ang{180} - \varpi$, $\cos v = -\cos(\lambda - \varpi)$.
Em $\lambda = \ang{90}$, $\cos(\lambda - \varpi) = \sin\varpi$ e $(a/r)^2 = (1 -
e\sin\varpi)^2/(1 - e^2)^2 = 1 - 2e\sin\varpi + O(e^2)$; em $\lambda = \ang{270}$ o sinal muda.
\end{proof}

A \omterm{def:b2:orbital-forcing:elements}{precessão climática} é portanto o índice certo: ela mede quanto mais perto do Sol, ou mais
longe, a Terra está no solstício de junho. Quando $e\sin\varpi$ é grande e positivo, como hoje e
muito mais há \num{116000} anos, o verão boreal cai perto do afélio e é fraco; quando é
grande e negativo, como há cerca de \num{10000} anos, o verão boreal cai perto do periélio.
% ledger: berger91:rows

\begin{example}[O Sol de junho a \ang{65} N, hoje]\label{ex:b2:orbital-forcing:today}
Hoje $e = 0.017236$, $\varpi = \ang{101.37}$, $\varepsilon = \ang{23.446}$, e as tabelas usam
$S_0 = \qty{1360}{W.m^{-2}}$. No solstício de junho $\delta = \varepsilon$; em $\varphi =
\ang{65}$, $\cos h_0 = -\tan\ang{65}\tan\ang{23.446} = -0.9301$, logo $h_0 = \ang{158.44} =
\qty{2.7654}{rad}$: o Sol fica acima do horizonte durante $2 \times 158.44/360 \times 24 = 21.1$~horas. O parêntese
da fórmula da insolação vale $2.7654 \times 0.9063 \times 0.3979 + 0.4226 \times 0.9174 \times 0.3674 =
1.1397$, e à distância média $\bar Q = 432.90 \times 1.1397 = \qty{493.4}{W.m^{-2}}$. O
fator de distância é $(a/r)^2 = 0.9671$, logo $\bar Q = \qty{477.1}{W.m^{-2}}$.
\end{example}
% ledger: berger91:rows, berger91:readme, data:b4-orbital

A segunda lei de Kepler acrescenta uma sutileza. Perto do afélio a Terra move-se mais devagar: a estação que cai
ali é mais fraca, mas dura mais. Os verões boreais são hoje cerca de 4.5 dias mais longos que os
invernos boreais, e as primaveras cerca de três dias mais longas que os outonos.
% ledger: nasa:milank-detail

\begin{proposition}[As médias anuais não dependem da precessão]\label{prop:b2:orbital-forcing:annual}
A energia recebida por unidade de área no topo da atmosfera entre duas longitudes dadas do
Sol não depende de $\varpi$, e a insolação média anual em qualquer latitude é a de uma
órbita circular de raio $a$ multiplicada por $1/\sqrt{1 - e^2}$. Num polo a média anual é
\[
  \bar Q_{\text{polo}} = \frac{S_0\sin\varepsilon}{\pi\sqrt{1 - e^2}} ,
\]
que cresce com a obliquidade.
\end{proposition}

\begin{proof}
Pela segunda lei de Kepler o raio varre área à taxa constante $\pi ab/T$, e a área varrida
enquanto a Terra gira de $\mathrm d\lambda$ é $\frac12 r^2\,\mathrm d\lambda$, logo $\mathrm dt =
T r^2\,\mathrm d\lambda/(2\pi ab)$, com $b = a\sqrt{1 - e^2}$. A insolação numa dada latitude
e longitude do Sol é $(a/r)^2$ vezes seu valor $Q_a(\lambda)$ à distância $a$, de modo que a
energia recebida é $\int Q_a(\lambda)(a/r)^2\,\mathrm dt = \frac{T}{2\pi\sqrt{1 - e^2}}\int
Q_a(\lambda)\,\mathrm d\lambda$: a distância cancelou-se, e com ela $\varpi$. Ao longo do ano
inteiro, $\lambda$ percorre \ang{360} e a média é $\frac{1}{2\pi\sqrt{1-e^2}}\int_0^{2\pi}
Q_a\,\mathrm d\lambda$. No Polo Norte $Q_a = S_0\sin\delta = S_0\sin\varepsilon\sin\lambda$
para $0 < \lambda < \pi$ e 0 no resto; a integral é $2S_0\sin\varepsilon$.
\end{proof}
% ledger: earth:TSI

A precessão desloca a luz solar de uma estação para outra, nunca para dentro ou para fora de um ano; a excentricidade
altera o total anual por um fator $1/\sqrt{1 - e^2}$, no máximo \qty{0.16}{\percent} no
último milhão de anos; a obliquidade desloca a luz solar entre latitudes, elevando a média anual nos
polos e baixando-a ligeiramente no equador. Só o padrão sazonal e latitudinal muda
muito, e é isso que a figura abaixo mostra.
% ledger: data:b4-orbital, nasa:milank-detail

\begin{omfigure}
\centering
% figdata/bachelor-2/orbital.py part lat
\begin{tikzpicture}
\begin{groupplot}[group style={group size=2 by 1, horizontal sep=1.6cm},
  tick label style={font=\scriptsize}, label style={font=\small}, height=5.4cm,
  width=0.48\linewidth, xmin=-90, xmax=90, xtick={-90,-60,...,90}, xlabel={latitude
  (\unit{\degree})}, ylabel={$\bar Q$ (\unit{W.m^{-2}})}]
\nextgroupplot[ymin=0, ymax=650, title={a\enspace solstício de junho}, title style={font=\scriptsize},
  legend style={font=\scriptsize, at={(0.03,0.97)}, anchor=north west, draw=none},
  legend cell align=left]
\addplot[black, thick] table[x=lat, y=jun_0] {figdata/out/bachelor-2-orbital-lat.dat};
\addlegendentry{hoje}
\addplot[ocRed, thick, densely dashed] table[x=lat, y=jun_10] {figdata/out/bachelor-2-orbital-lat.dat};
\addlegendentry{há \num{10000} anos}
\nextgroupplot[ymin=100, ymax=500, title={b\enspace média anual}, title style={font=\scriptsize},
  legend style={font=\scriptsize, at={(0.5,0.03)}, anchor=south, draw=none},
  legend cell align=left]
\addplot[ocBlue, thick] table[x=lat, y=ann_lo] {figdata/out/bachelor-2-orbital-lat.dat};
\addlegendentry{$\varepsilon = \ang{22.1}$}
\addplot[ocRed, thick, densely dashed] table[x=lat, y=ann_hi] {figdata/out/bachelor-2-orbital-lat.dat};
\addlegendentry{$\varepsilon = \ang{24.5}$}
\end{groupplot}
\end{tikzpicture}
\par\smallskip{\small (a) Insolação média diária no solstício de junho, hoje e há \num{10000}
anos, quando o verão boreal caía perto do periélio e a obliquidade era maior: cerca de
\qty{50}{W.m^{-2}} a mais nas altas latitudes do norte. (b) Insolação média anual para as duas
obliquidades extremas do último milhão de anos, com a excentricidade atual: uma inclinação maior manda
mais luz solar para os polos e um pouco menos para os trópicos. Calculadas a partir dos elementos orbitais
com $S_0 = \qty{1360}{W.m^{-2}}$.}
\end{omfigure}
% ledger: data:b4-orbital, data:b4-orbit91

\section{Calcular a curva de insolação}\label{sec:b2:orbital-forcing:curve}

A curva de insolação que Milankovitch construiu é hoje um cálculo de rotina. Os elementos orbitais do
passado são calculados integrando os movimentos dos planetas; tabelados a cada mil anos,
eles são os dados de entrada. Para cada idade, o método abaixo dá a insolação em qualquer latitude e estação.

\begin{method}[A insolação do passado]\label{met:b2:orbital-forcing:curve}
\begin{enumerate}
  \item Ler $e$, $\varpi$ e $\varepsilon$ para a idade desejada.
  \item Escolher a estação pela longitude verdadeira do Sol $\lambda$: \ang{90} para o solstício de junho;
        o meio de julho das tabelas corresponde a $\lambda = \ang{120}$.
  \item Calcular $\delta$ a partir de $\sin\delta = \sin\varepsilon\sin\lambda$, e o ângulo horário do
        pôr do sol $h_0$ na latitude $\varphi$ (dia ou noite polar se $|\tan\varphi\tan\delta| > 1$).
  \item Calcular o fator de distância $(a/r)^2$ com a \cref{prop:b2:orbital-forcing:distance}.
  \item Multiplicar: $\bar Q = (S_0/\pi)(a/r)^2(h_0\sin\varphi\sin\delta +
        \cos\varphi\cos\delta\sin h_0)$.
\end{enumerate}
\end{method}
% ledger: berger91:readme, data:b4-orbital

Aplicado às 1001 idades do último milhão de anos, o método reproduz a insolação tabelada de meados de julho
a \ang{65} N com diferença de até \qty{0.012}{W.m^{-2}}, o arredondamento das tabelas, e os
valores tabelados de junho no presente exatamente. As tabelas põem a origem do tempo em 1950.
% ledger: data:b4-orbital, berger91:readme

\begin{omfigure}
\centering
% figdata/bachelor-2/orbital.py part jul65
\begin{tikzpicture}
\begin{groupplot}[group style={group size=1 by 2, vertical sep=0.3cm, xlabels at=edge bottom,
  xticklabels at=edge bottom}, width=0.96\linewidth, xmin=0, xmax=400, xtick={0,50,...,400},
  tick label style={font=\scriptsize}, label style={font=\small}]
\nextgroupplot[height=5.0cm, ymin=370, ymax=590, ylabel={$\bar Q$ a \ang{65} N
  (\unit{W.m^{-2}})}, legend style={font=\scriptsize, at={(0.5,1.02)}, anchor=south, draw=none,
  legend columns=3, /tikz/every even column/.append style={column sep=6pt}},
  legend cell align=left]
\addplot[black!45, line width=2.2pt] table[x=age_ka, y=tab] {figdata/out/bachelor-2-orbital-jul65.dat};
\addlegendentry{meados de julho, tabelado}
\addplot[ocRed, thin] table[x=age_ka, y=calc] {figdata/out/bachelor-2-orbital-jul65.dat};
\addlegendentry{meados de julho, calculado}
\addplot[ocBlue, thin, densely dashed] table[x=age_ka, y=jun65] {figdata/out/bachelor-2-orbital-jul65.dat};
\addlegendentry{solstício de junho, calculado}
\nextgroupplot[height=2.8cm, ymin=-0.02, ymax=0.02, ytick={-0.01,0,0.01},
  scaled y ticks=false, yticklabel style={/pgf/number format/fixed}, ylabel={resíduo},
  xlabel={idade (milhares de anos antes de 1950)}]
\addplot[only marks, mark=*, mark size=0.35pt, black] table[x=age_ka, y=resid]
  {figdata/out/bachelor-2-orbital-jul65.dat};
\end{groupplot}
\end{tikzpicture}
\par\smallskip{\small Em cima: insolação média diária a \ang{65} N nos últimos \num{400000} anos,
em meados de julho tal como tabelada por Berger e Loutre (cinza, largo) e tal como calculada pelo método acima (vermelho),
e no solstício de junho (azul, tracejado). Embaixo: valores de meados de julho calculados menos tabelados, em
\unit{W.m^{-2}}.}
\end{omfigure}
% ledger: data:b4-orbit91, data:b4-orbital

Dois traços se destacam. A curva é dominada pela \omterm{def:b2:orbital-forcing:elements}{precessão climática}, com picos separados por cerca de
\num{20000} anos cuja altura aumenta e diminui com a excentricidade; a obliquidade acrescenta
uma modulação mais lenta. E as oscilações são grandes: nos últimos \num{800000} anos o valor de meados de julho
variou de 388 a \qty{497}{W.m^{-2}}.
% ledger: data:b4-orbital

\section{A teoria de Milankovitch e os ciclos glaciais}\label{sec:b2:orbital-forcing:theory}

\begin{history}[A insolação a sessenta e cinco graus norte]
Milutin Milankovitch, nascido em 1879, construiu uma teoria matemática do clima a partir das variações sazonais e latitudinais da luz solar. Combinando as três variações orbitais,
calculou a insolação das latitudes e as temperaturas de superfície correspondentes para
\num{600000} anos antes de 1800. Por sugestão do climatólogo Köppen, escolheu a insolação de verão
a \ang{65} N como a grandeza decisiva: os grandes mantos de gelo tinham crescido perto dessa
latitude, e os verões frescos derretem menos da neve do inverno, de modo que a neve se acumula de ano para
ano. Ele considerava a obliquidade o mais importante dos três parâmetros. Durante cerca de cinquenta anos
a teoria foi em grande parte ignorada; Milankovitch morreu em 1958. Em 1976 Hays, Imbrie e Shackleton
encontraram em sedimentos de mar profundo dos últimos \num{450000} anos picos espectrais de variabilidade climática
em cerca de \num{23000}, \num{42000} e \num{100000} anos, os períodos da precessão, da obliquidade e da
excentricidade, e concluíram que as variações da órbita terrestre são a causa fundamental das
eras glaciais quaternárias.
\end{history}
% ledger: eo00:milankovitch, nasa:milank-detail, hays76:periods

A figura abaixo repete seu teste com os registros do \cref{ch:b2:climate-proxies}. Ela compara os
espectros de amplitude, a intensidade de cada período, da insolação de julho a \ang{65} N e da
compilação bentônica, em duas janelas de \num{700000} anos.

\begin{omfigure}
\centering
% figdata/bachelor-2/orbital.py part spectrum
\begin{tikzpicture}
\begin{axis}[width=0.9\linewidth, height=5.6cm, xmin=10, xmax=150, ymin=0, ymax=1.08,
  xlabel={período (milhares de anos)}, ylabel={amplitude relativa},
  xtick={19,23,41,100,150}, tick label style={font=\scriptsize}, label style={font=\small},
  legend style={font=\scriptsize, at={(0.5,1.02)}, anchor=south, draw=none, legend columns=3,
  /tikz/every even column/.append style={column sep=6pt}}, legend cell align=left]
\addplot[black!50, thick] table[x=period, y=amp_ins] {figdata/out/bachelor-2-orbital-spectrum.dat};
\addlegendentry{insolação \ang{65} N, 0--\qty{700}{ka}}
\addplot[ocRed, thick] table[x=period, y=amp_young] {figdata/out/bachelor-2-orbital-spectrum.dat};
\addlegendentry{LR04, 0--\qty{700}{ka}}
\addplot[ocBlue, thick, densely dashed] table[x=period, y=amp_old] {figdata/out/bachelor-2-orbital-spectrum.dat};
\addlegendentry{LR04, de 1250 a 1950~ka}
\end{axis}
\end{tikzpicture}
\par\smallskip{\small Espectros de amplitude, cada um normalizado pelo seu próprio máximo, da insolação de meados de julho
a \ang{65} N nos últimos \num{700000} anos e da compilação bentônica LR04 na mesma janela
e de \num{1250000} a \num{1950000} anos atrás (tendência removida, com janela de ponderação). A insolação tem sua
potência nos períodos da precessão, perto de 19 a 24 mil anos, e no período da obliquidade; o volume de gelo
recente tem seu pico perto de \num{100000} anos, onde a insolação quase nada tem; o volume de gelo mais antigo
tem seu pico em \num{41000} anos.}
\end{omfigure}
% ledger: data:b4-orbital, data:lr04, hays76:periods

\subsection*{Glaciais, interglaciais e terminações}

\begin{definition}[Período glacial; interglacial; terminação glacial]\label{def:b2:orbital-forcing:glacials}
No registro quaternário, um \emph{período glacial}\index{período glacial} é um intervalo em que grandes
mantos de gelo cobrem os continentes do norte, registrado por um $\delta^{18}$O bentônico alto e um nível do mar baixo;
um \emph{interglacial}\index{interglacial} é um intervalo com pouco gelo continental no Hemisfério
Norte fora da Groenlândia, como o presente. Uma \emph{terminação
glacial}\index{terminação glacial} é uma das grandes transições que encerram um ciclo glacial,
de condições plenamente glaciais a condições interglaciais, muito mais rápida que o crescimento que a precedeu.
\end{definition}
% ledger: kohler20:mpt, barker21:terminations, ipcc:ref-periods

No último meio milhão de anos o período dominante do clima é próximo de \num{100000} anos,
embora o forçamento nesse período, que vem só da excentricidade, seja desprezível
diante do da obliquidade e da precessão. Os ciclos são também assimétricos: o $\delta^{18}$O bentônico
registra uma acumulação de gelo em degraus ao longo de cerca de \num{90000} anos, depois um recuo em cerca de
\num{10000}, um dente de serra cujas bordas descendentes são as terminações. O último \omterm{def:b2:orbital-forcing:glacials}{interglacial}, de
cerca de \num{129000} a \num{116000} anos atrás, terminou quando a insolação de verão a \ang{65} N caiu
para um de seus valores mais baixos dos últimos \num{800000} anos: o valor do solstício de junho atingiu seu
mínimo, \qty{438.6}{W.m^{-2}}, há \num{116000} anos, e o valor de meados de julho seu segundo mínimo mais baixo,
\qty{392.2}{W.m^{-2}}, há \num{114000} anos. O Último Máximo Glacial veio cerca de
\num{95000} anos depois.
% ledger: barker21:terminations, ipcc:ref-periods, data:b4-orbital, data:b4-orbital-extremes

Com que rapidez um manto de gelo pode seguir seu forçamento? O modelo mais simples trata o volume de gelo $V$ como
relaxando em direção ao volume $V_{\mathrm{eq}}$ que a insolação presente sustentaria, com um
tempo de resposta $\tau$:
\[
  \frac{\mathrm dV}{\mathrm dt} = \frac{V_{\mathrm{eq}}(t) - V}{\tau} .
\]

\begin{proposition}[Uma resposta defasada]\label{prop:b2:orbital-forcing:lag}
Para um forçamento $V_{\mathrm{eq}} = A\sin\omega t$ a solução do modelo de atraso tende para
\[
  V = \frac{A}{\sqrt{1 + \omega^2\tau^2}}\,\sin\bigl(\omega(t - t_{\ell})\bigr),\qquad
  t_{\ell} = \frac{\arctan(\omega\tau)}{\omega} :
\]
o gelo segue quase totalmente um forçamento lento e só fracamente um forçamento rápido, sempre com atraso. A resposta
a uma soma de senoides é a soma das respostas: o modelo não cria nenhum período que falte a seu
forçamento.
\end{proposition}

\begin{proof}
Tente-se $V = B\sin\omega t + C\cos\omega t$: substituindo, $\omega B\tau = -C$ e $-\omega C\tau = A -
B$, logo $B = A/(1 + \omega^2\tau^2)$, $C = -\omega\tau A/(1 + \omega^2\tau^2)$, uma amplitude
$\sqrt{B^2 + C^2} = A/\sqrt{1 + \omega^2\tau^2}$ e um atraso de fase $\arctan(\omega\tau)$. Qualquer outra
solução difere desta por um múltiplo de $\mathrm e^{-t/\tau}$, que se extingue. A
equação é linear, de modo que as soluções se somam.
\end{proof}

\begin{omfigure}
\centering
% figdata/bachelor-2/orbital.py part lag
\begin{tikzpicture}
\begin{axis}[width=0.96\linewidth, height=5.4cm, xmin=0, xmax=400, ymin=-3.4, ymax=4.2,
  xtick={0,50,...,400}, xlabel={idade (milhares de anos antes de 1950)},
  ylabel={padronizado}, tick label style={font=\scriptsize}, label style={font=\small},
  legend style={font=\scriptsize, at={(0.5,1.02)}, anchor=south, draw=none, legend columns=3,
  /tikz/every even column/.append style={column sep=6pt}}, legend cell align=left]
\addplot[black!35, thin] table[x=age_ka, y=forcing] {figdata/out/bachelor-2-orbital-lag.dat};
\addlegendentry{menos a insolação \ang{65} N}
\addplot[ocRed, thick] table[x=age_ka, y=model] {figdata/out/bachelor-2-orbital-lag.dat};
\addlegendentry{modelo de atraso, $\tau = \qty{10}{kyr}$}
\addplot[black, thick] table[x=age_ka, y=lr04] {figdata/out/bachelor-2-orbital-lag.dat};
\addlegendentry{LR04 (gelo para cima)}
\end{axis}
\end{tikzpicture}
\par\smallskip{\small O modelo de atraso movido pela insolação de meados de julho a \ang{65} N (menos
sol, mais gelo), comparado com a compilação bentônica, as três padronizadas (média 0, desvio-padrão
1); os valores positivos significam mais gelo. O modelo segue as ondulações de precessão da compilação, mas não tem
nenhum de seus grandes ciclos em dente de serra. O tempo de resposta de \qty{10}{kyr} é uma escolha do modelo.}
\end{omfigure}
% ledger: data:b4-orbital, data:lr04

Nos últimos \num{800000} anos, o modelo defasado correlaciona-se com a compilação com $r = 0.42$, a
própria insolação com apenas $0.23$: o atraso ajuda, mas um manto de gelo linear não consegue fazer
ciclos de \num{100000} anos a partir de um forçamento que não os tem. As terminações devem envolver
um comportamento não linear que amplifica um empurrão modesto da insolação. Os candidatos são retroalimentações que agem sobre o clima inteiro:
a \omterm{def:b2:precambrian-earth:ice-albedo}{retroalimentação gelo--albedo} da \cref{def:b2:precambrian-earth:ice-albedo}, já que um planeta com mais
neve e gelo reflete mais luz solar; e o dióxido de carbono, que nos testemunhos de gelo antárticos sobe em
escalas de milênios sempre que a circulação de revolvimento do Atlântico enfraquece, como acontece
quando a água de degelo dos mantos de gelo se despeja no Atlântico Norte. Acima de um limiar de dióxido
de carbono, as condições \omterm{def:b2:orbital-forcing:glacials}{interglaciais} podem estabilizar-se.
% ledger: data:b4-orbital, barker21:terminations, eo00:milankovitch

\begin{definition}[Transição do Pleistoceno Médio]\label{def:b2:orbital-forcing:mpt}
A \emph{transição do Pleistoceno Médio}\index{transição do Pleistoceno Médio} é a mudança, de cerca de
\num{1250000} a \num{700000} anos atrás, de ciclos glaciais de cerca de \num{41000} anos e pequena
amplitude para ciclos quase periódicos de cerca de \num{100000} anos e grande amplitude, com mantos de gelo
mais espessos, sem nenhuma mudança correspondente no forçamento orbital.
\end{definition}
% ledger: clark06:mpt, nasa:milank-detail, kohler20:mpt

Antes da transição os mantos de gelo seguiam a obliquidade, como mostra o espectro tracejado da figura
acima; uma contagem encontra uma \omterm{def:b2:orbital-forcing:glacials}{terminação glacial} em 67, 88 e \qty{52}{\percent} dos
ciclos de obliquidade do Quaternário inicial, médio e final. A teoria orbital sozinha não explica a mudança; os candidatos são a dinâmica
dos próprios mantos de gelo, talvez depois da remoção do regolito mole sobre o qual deslizavam,
e um declínio gradual do dióxido de carbono durante os mínimos glaciais. Nenhuma explicação é ainda
geralmente aceita.
% ledger: kohler20:mpt, nasa:milank-detail, clark06:mpt

\section{O dióxido de carbono nos ciclos glaciais}\label{sec:b2:orbital-forcing:co2}

Os testemunhos de gelo antárticos mostram que o dióxido de carbono acompanha os ciclos glaciais: entre 188 e
\qty{194}{ppm} no Último Máximo Glacial, entre 266 e \qty{282}{ppm} no último \omterm{def:b2:orbital-forcing:glacials}{interglacial};
o registro composto do \cref{ch:b2:climate-proxies} dá \qty{189.7}{ppm} para 19 a
\qty{23}{ka} e \qty{279.4}{ppm} nos dois mil anos anteriores à indústria. Seu forçamento radiativo
cresce perto do logaritmo da concentração, \qty{3.93}{W.m^{-2}} por duplicação.
% ledger: ipcc:fig2-34, data:b4-proxies-co2, ipcc:log-co2, ipcc:erf-2xco2

\begin{example}[A órbita contra o dióxido de carbono]\label{ex:b2:orbital-forcing:co2}
De \qty{189.7}{ppm} a \qty{279.4}{ppm}, o forçamento é $3.93\log_2(279.4/189.7) = 3.93 \times
0.559 = \qty{2.2}{W.m^{-2}}$, sobre o globo inteiro e o ano inteiro. A órbita, no último
milhão de anos, alterou a média global e anual da insolação em no máximo $340.54 - 340.00 =
\qty{0.54}{W.m^{-2}}$, mas a insolação de meados de julho a \ang{65} N em cerca de \qty{109}{W.m^{-2}}.
\end{example}
% ledger: data:b4-proxies-co2, ipcc:erf-2xco2, ipcc:log-co2, data:b4-orbital

A órbita é um marca-passo, não um motor. Ela redistribui a luz solar entre estações e latitudes
e fixa o calendário; os mantos de gelo, por seu albedo e sua lentidão, e o dióxido de carbono, por
seu forçamento global, transformam um sinal regional numa mudança global do clima.
% ledger: hays76:periods, barker21:terminations, eo00:milankovitch

\section*{Exercícios}\label{sec:b2:orbital-forcing:exercises}

Dados para os exercícios, salvo indicação em contrário: $S_0 = \qty{1360}{W.m^{-2}}$ (o valor das
tabelas); um ano trópico de \num{365.2422} dias; os elementos orbitais das tabelas, dos quais os de hoje são
$e = 0.017236$, $\varpi = \ang{101.37}$, $\varepsilon = \ang{23.446}$.
% ledger: berger91:readme, berger91:rows

\begin{exercise}[$\star$]\label{exo:b2:orbital-forcing:1}
Tome $e = 0.0167$. (a) Calcule a razão entre a distância no afélio $a(1 + e)$ e a distância no
periélio $a(1 - e)$. (b) Exprima sua diferença em porcentagem de $a$. (c) Por que fator
a luz solar no periélio é mais forte que no afélio?
\end{exercise}
% ledger: earth:ecc, nasa:milank-detail

\begin{exercise}[$\star$]\label{exo:b2:orbital-forcing:2}
As tabelas dão $e = 0.017236$, $\varpi = \ang{101.37}$ hoje; $e = 0.019971$, $\varpi =
\ang{293.86}$ há \num{10000} anos; $e = 0.019249$, $\varpi = \ang{359.99}$ há \num{6000} anos.
(a) Calcule a \omterm{def:b2:orbital-forcing:elements}{precessão climática} para cada idade. (b) Para cada uma, dê a longitude do Sol no
periélio. (c) Em que idade o periélio estava mais perto do solstício de junho, e em qual do
equinócio de setembro?
\end{exercise}
% ledger: berger91:rows

\begin{exercise}[$\star$]\label{exo:b2:orbital-forcing:3}
O eixo de rotação precessiona com um período de cerca de \num{25770} anos em relação às estrelas, enquanto
a órbita gira com um período de cerca de \num{112000} anos. Suponha que os dois movimentos desloquem o
periélio em relação ao equinócio no mesmo sentido, de modo que suas frequências se somem. (a) Combine
as duas frequências para encontrar o período com que o periélio
percorre as estações. (b) Compare com os períodos principais da \omterm{def:b2:orbital-forcing:elements}{precessão climática}, \num{19.0},
\num{22.4} e \num{23.7} mil anos.
\end{exercise}
% ledger: nasa:milank-detail, la04:precession

\begin{exercise}[$\star$]\label{exo:b2:orbital-forcing:4}
Com a \cref{prop:b2:orbital-forcing:annual} e $e = 0.0167$, calcule a insolação média anual no
Polo Norte para obliquidades de \ang{22.1} e \ang{24.5}. Em que porcentagem ela cresce entre
as duas?
\end{exercise}
% ledger: nasa:milank-detail, earth:ecc

\begin{exercise}[$\star\star$]\label{exo:b2:orbital-forcing:5}
(a) Calcule exatamente o fator de distância $(a/r)^2$ no solstício de junho hoje e há \num{116000} anos
($e = 0.043988$, $\varpi = \ang{92.71}$), e compare com a aproximação $1 -
2e\sin\varpi$. (b) Calcule ambos no solstício de dezembro. (c) Por que a aproximação é pior há \num{116000}
anos?
\end{exercise}
% ledger: berger91:rows

\begin{exercise}[$\star\star$]\label{exo:b2:orbital-forcing:6}
Reproduza o \cref{ex:b2:orbital-forcing:today} passo a passo: o ângulo horário do pôr do sol a \ang{65} N
no solstício de junho, a duração do dia, a insolação à distância média e a
insolação de hoje.
\end{exercise}
% ledger: berger91:rows

\begin{exercise}[$\star\star$]\label{exo:b2:orbital-forcing:7}
(a) Usando $\mathrm dt = T(r/a)^2\,\mathrm d\lambda/(2\pi\sqrt{1 - e^2})$ e $(r/a)^2 \approx 1 +
2e\cos(\lambda - \varpi)$, mostre que o semestre do equinócio de março ao de setembro dura
cerca de $T/2 + (2T/\pi)\,e\sin\varpi$. (b) Calcule em quantos dias ele excede o outro semestre
hoje. (c) Compare com os 4.5 dias em que os verões boreais excedem os invernos e os três dias em
que as primaveras excedem os outonos.
\end{exercise}
% ledger: nasa:milank-detail, berger91:rows

\begin{exercise}[$\star\star$]\label{exo:b2:orbital-forcing:8}
No último milhão de anos a excentricidade variou de \num{0.0046} a \num{0.0565}. (a) Calcule
a insolação média global e anual $S_0/(4\sqrt{1 - e^2})$ nos dois extremos. (b) Em quanto, em
\unit{W.m^{-2}} e em porcentagem, a órbita a alterou? (c) Compare com a amplitude da
insolação de meados de julho a \ang{65} N, 388 a \qty{497}{W.m^{-2}}.
\end{exercise}
% ledger: data:b4-orbital

\begin{exercise}[$\star\star$]\label{exo:b2:orbital-forcing:9}
Com o modelo de atraso e $\tau = \qty{10}{kyr}$, calcule a amplitude relativa e o atraso do
volume de gelo para forçamentos senoidais de períodos 23, 41 e \num{100} mil anos. Que banda
o manto de gelo favorece, e por quê?
\end{exercise}

\begin{exercise}[$\star\star\star$]\label{exo:b2:orbital-forcing:10}
A insolação de meados de julho a \ang{65} N quase não tem potência em \num{100000} anos, e no entanto o volume de gelo
dos últimos \num{700000} anos tem ali a maior parte da sua. (a) Mostre que, se o forçamento é $A_1\sin
\omega_1 t + A_2\sin\omega_2 t$, o modelo de atraso responde apenas em $\omega_1$ e $\omega_2$. (b) Mostre
que uma resposta que contém o termo $V^2$, por exemplo $\mathrm dV/\mathrm dt = (V_{\mathrm{eq}} -
V)/\tau + kV^2$, pode criar as frequências $\omega_1 - \omega_2$ e $\omega_1 + \omega_2$,
substituindo a solução linear em $V^2$. (c) Dois termos de precessão de períodos \num{23.7} e
\num{19.0} mil anos batem com que período? O que você conclui sobre os ciclos de
\num{100000} anos?
\end{exercise}
% ledger: la04:precession, barker21:terminations

\begin{exercise}[$\star\star\star$]\label{exo:b2:orbital-forcing:11}
No Polo Norte, a energia recebida por metro quadrado durante o dia polar é $E = \int S_0(a/r)^2
\sin\delta\,\mathrm dt$ ao longo do semestre em que $\delta > 0$. (a) Mostre com a demonstração da
\cref{prop:b2:orbital-forcing:annual} que $E = S_0T\sin\varepsilon/(\pi\sqrt{1 - e^2})$. (b) Calcule
$E$ em \unit{GJ.m^{-2}} para $\varepsilon = \ang{22.1}$ e \ang{24.5}, com $e = 0.0167$ e $T$ =
365.25 dias. (c) A precessão não aparece em $E$. Por que ela pode, mesmo assim, alterar o
derretimento de um manto de gelo?
\end{exercise}
% ledger: nasa:milank-detail

\begin{exercise}[$\star\star\star$]\label{exo:b2:orbital-forcing:12}
(a) Calcule o forçamento radiativo da subida do dióxido de carbono de \qty{189.7}{ppm} para
\qty{279.4}{ppm}, com \qty{3.93}{W.m^{-2}} por duplicação e uma lei logarítmica. (b) Os intervalos
dados para o Último Máximo Glacial e o último \omterm{def:b2:orbital-forcing:glacials}{interglacial} são 188 a \qty{194}{ppm} e 266 a
\qty{282}{ppm}: calcule o menor e o maior forçamento que eles permitem. (c) Compare com a variação
da insolação global anual do \cref{exo:b2:orbital-forcing:8} e com a oscilação da insolação de verão a
\ang{65} N. Qual é o gatilho e qual o amplificador?
\end{exercise}
% ledger: data:b4-proxies-co2, ipcc:fig2-34, ipcc:erf-2xco2, ipcc:log-co2, data:b4-orbital

\begin{problem}[{Problema de fim de semana --- a luz solar a \ang{65} N}]\label{pb:b2:orbital-forcing:1}
O último \omterm{def:b2:orbital-forcing:glacials}{interglacial} terminou há cerca de \num{116000} anos. Quanta luz solar recebia então o verão da
latitude dos grandes mantos de gelo do norte, em comparação com hoje? Use os dados dos
exercícios. As tabelas dão, para \num{116000} anos atrás, $e = 0.0440$, $\varpi = \ang{92.7}$,
$\varepsilon = \ang{22.52}$; hoje, arredondados da mesma maneira, $e = 0.0172$, $\varpi = \ang{101.4}$,
$\varepsilon = \ang{23.45}$.

\medskip\noindent\textbf{Parte I --- A órbita então e agora.}
\begin{enumerate}
  \item Calcule a \omterm{def:b2:orbital-forcing:elements}{precessão climática} então e agora.
  \item Calcule as distâncias no periélio e no afélio, em unidades de $a$, então e agora.
  \item Em que longitude do Sol a Terra estava no periélio então, e agora? Perto de que solstício?
  \item Por que fator a luz solar no periélio era mais forte que no afélio, então e agora?
  \item A obliquidade era maior ou menor que hoje? O que isso significa para os verões das
        altas latitudes?
  \item Calcule a insolação média global e anual então e agora. A diferença é significativa?
\end{enumerate}

\medskip\noindent\textbf{Parte II --- A distância no solstício de junho.}
\begin{enumerate}[resume]
  \item Mostre a partir da \cref{prop:b2:orbital-forcing:distance} que no solstício de junho $r/a = (1 -
        e^2)/(1 - e\sin\varpi)$.
  \item Calcule $r/a$ no solstício de junho então e agora.
  \item Calcule $(a/r)^2$ então e agora.
  \item Compare com a aproximação $1 - 2e\sin\varpi$.
  \item Em que porcentagem a luz solar do solstício era mais fraca que à distância média, então e agora?
\end{enumerate}

\medskip\noindent\textbf{Parte III --- A insolação a \ang{65} N.}
\begin{enumerate}[resume]
  \item Dê a declinação do Sol no solstício de junho então e agora.
  \item Calcule o ângulo horário do pôr do sol a \ang{65} N nesse dia, então e agora.
  \item Deduza a duração do dia.
  \item Calcule a insolação média diária à distância média.
  \item Calcule a insolação média diária, então e agora.
  \item As tabelas dão, para meados de julho, 397.60 e \qty{426.76}{W.m^{-2}}. Por que os valores do
        solstício são mais altos?
  \item O que pesou mais na diferença entre então e agora, a obliquidade ou a distância?
\end{enumerate}

\medskip\noindent\textbf{Parte IV --- A resposta do manto de gelo.}
\begin{enumerate}[resume]
  \item No modelo de atraso com $\tau = \qty{10}{kyr}$, que fração de uma variação súbita de
        $V_{\mathrm{eq}}$ o manto de gelo realizou depois de 5, 10 e \qty{20}{kyr}?
  \item A insolação de junho a \ang{65} N caiu de um máximo há cerca de \num{128000} anos até o
        mínimo de \num{116000} anos atrás, meio ciclo de precessão. Para uma senoide de período
        \qty{23}{kyr}, que fração do forçamento o manto de gelo segue, e com que atraso?
  \item Por que um manto de gelo não pode atingir, num só ciclo de precessão, o tamanho que uma insolação de verão
        baixa sustentaria?
  \item O gelo levou uns \num{90000} anos para se acumular até o Último Máximo Glacial e cerca de
        \num{10000} anos para derreter. Por que o modelo de atraso não pode produzir essa assimetria?
  \item O dióxido de carbono estava entre 266 e \qty{282}{ppm} no último \omterm{def:b2:orbital-forcing:glacials}{interglacial} e entre
        188 e \qty{194}{ppm} no Último Máximo Glacial. Como ele ajuda uma mudança regional de
        insolação a tornar-se uma mudança global do clima?
  \item Em resumo: qual era a insolação média diária a \ang{65} N no solstício de junho quando o
        último \omterm{def:b2:orbital-forcing:glacials}{interglacial} terminou, há \num{116000} anos?
\end{enumerate}
\end{problem}
% ledger: berger91:rows, berger91:readme, data:b4-orbital, ipcc:ref-periods, ipcc:fig2-34, barker21:terminations
